\documentclass[12pt,a4paper]{article}
\usepackage[utf8]{inputenc}
\usepackage[vietnamese]{babel}
\usepackage[T5]{fontenc}
\usepackage{amsmath,amssymb,amsfonts}
\usepackage{fancyhdr}
\usepackage{geometry}
\usepackage{tcolorbox}

\geometry{
	a4paper,
	left=20mm,
	right=20mm,
	top=20mm,
	bottom=20mm
}

\pagestyle{fancy}
\fancyhf{}
\lhead{\small\textit{Bổ đề Hình học: Tứ giác điều hòa, Tam giác trực tâm \& Phép vị tự}}
\rhead{\small\textbf{Integra Typesetter}}
\rfoot{\small Trang \thepage}
\renewcommand{\headrulewidth}{0.4pt}
\renewcommand{\footrulewidth}{0.4pt}

\begin{document}

\begin{center}
	{\Large \textbf{CÁC BỔ ĐỀ HÌNH HỌC PHẲNG NÂNG CAO}}\\[2mm]
	{\small \textit{(Tứ giác điều hòa -- Định lý Ptolemy -- Tam giác trực tâm \& Phép vị tự)}}
\end{center}
\vspace{4mm}

\begin{tcolorbox}[colback=blue!5!white,colframe=blue!75!black,title=\textbf{Bổ đề 1: Tứ giác điều hòa và Định lý Ptolemy}]
Cho đường tròn $(I)$ và điểm $A$ nằm ngoài đường tròn với hai tiếp tuyến $AN, AP$. Cát tuyến bất kỳ qua $A$ cắt $(I)$ tại $A'$ và $D$. Khi đó:
\begin{enumerate}
    \item Tứ giác $NPA'D$ là tứ giác điều hòa, thỏa mãn:
    \[
    DN \cdot PA' = DP \cdot NA' = \frac{1}{2} DA' \cdot NP
    \]
    \item Hệ thức Ptolemy: $DN \cdot PA' + DP \cdot NA' = DA' \cdot NP$.
\end{enumerate}
\end{tcolorbox}

\vspace{3mm}

\begin{tcolorbox}[colback=teal!5!white,colframe=teal!75!black,title=\textbf{Bổ đề 2: Tâm nội tiếp và Ngoại tiếp của Tam giác trực tâm}]
Cho tam giác nhọn $ABC$ có trực tâm $H$, tâm ngoại tiếp $O$ và tam giác trực tâm là $UVW$ (chân các đường cao). Khi đó:
\begin{enumerate}
    \item $H$ là tâm đường tròn nội tiếp của tam giác $UVW$.
    \item Tâm đường tròn ngoại tiếp của tam giác $UVW$ là tâm đường tròn Euler $O_9$ của tam giác $ABC$.
    \item Đường thẳng nối tâm nội tiếp và tâm ngoại tiếp của tam giác $UVW$ trùng với đường thẳng Euler $OH$ của tam giác $ABC$.
\end{enumerate}
\end{tcolorbox}

\vspace{3mm}

\begin{tcolorbox}[colback=orange!5!white,colframe=orange!75!black,title=\textbf{Bổ đề 3: Phép vị tự và Trục nối các tâm đặc biệt}]
Nếu hai tam giác $T_1$ và $T_2$ vị tự với nhau (có các cạnh tương ứng song song), thì đường thẳng nối tâm nội tiếp và tâm ngoại tiếp của $T_1$ song song (hoặc trùng) với đường thẳng nối tâm nội tiếp và tâm ngoại tiếp của $T_2$.
\end{tcolorbox}

\end{document}
